Now $H$ is the spectrum of a finite local $k$-algebra with residue field $k$ (cf.\ VI A, 5.6.1). Consequently, there exists an integer $n_0$ such that, for every $n\geq n_0$, $\mathrm{Fr}^n(H/k)$ factors through the ``unit'' section of $H^{(p^n)}$. It follows that, for $n\geq n_0$, $\mathrm{Fr}^n(G/k)$ factors through $G_{\mathrm{red}}^{(p^n)}$, and hence, according to VI A, 5.4.1, there is a commutative diagram
\[
\xymatrix@C=4em{G\ar[r]^{\mathrm{Fr}^n(G/k)}\ar[d]_\pi&G_{\mathrm{red}}^{(p^n)}\\G/({}_{\mathrm{Fr}^n}G)\ar[ur]_i&}
\]
where $i$ is a closed immersion (and $\pi$ the canonical projection). Since, in addition, $i$ induces a homeomorphism of the underlying topological spaces, it is therefore an isomorphism. Since $k$ is perfect, $G_{\mathrm{red}}^{(p^n)}$ is smooth over $k$ (VI A, 1.3.1), and hence $G/({}_{\mathrm{Fr}^n}G)$ is smooth over $k$, for every $n\geq n_0$.

\begin{corollary}{8.3.1}\label{VIIA.8.3.1}
Let $G$ be a group locally of finite type\footnotemark[80] over a field $k$ of characteristic $p>0$, and let $n$ be an integer $\geq1$.\nde{The equivalence between (ii) and (iii) has been made explicit, and the proof expanded.} The following conditions are equivalent:
\begin{enumeratei}
\item $G$ is smooth over $k$.
\item $\mathrm{Fr}^n(G/k):G\to G^{(p^n)}$ is an epimorphism of (fppf) sheaves.
\item $\mathrm{Fr}^n(G/k):G\to G^{(p^n)}$ is faithfully flat.
\end{enumeratei}
\end{corollary}
First, since $G$ is locally of finite type over $k$, $\mathrm{Fr}^n(G/k)$ is of finite presentation, so the equivalence of (ii) and (iii) follows from 8.1.4. Suppose $G$ smooth over $k$, hence $G$ reduced. Then, since $\mathrm{Fr}^n(G/k)$ is surjective, it is faithfully flat (cf.\ VI A, 6.2 or VI B, 1.3).

Conversely, suppose $\mathrm{Fr}^n(G/k)$ faithfully flat. Since $\mathrm{Fr}^n(G^{(p^n)}/k)$ is deduced from $\mathrm{Fr}^n(G/k)$ by base change (cf.\ 4.1.3), it too is faithfully flat, as is the composite:
\[
\mathrm{Fr}^{2n}(G/k):G\longrightarrow G^{(p^n)}\longrightarrow G^{(p^{2n})}.
\]
One thus obtains that, for every $m\in\NN$, $\mathrm{Fr}^{mn}(G/k):G\to G^{(p^{mn})}$ is faithfully flat, hence induces an isomorphism $G/({}_{\mathrm{Fr}^{mn}}G)\simeq G^{(p^{mn})}$ (cf.\ VI A, 5.4.1). Now, according to proposition 8.3, $G^{(p^{mn})}$ is smooth over $k$ for large $m$, hence so is $G$, by (fpqc) descent (cf.\ EGA~IV$_4$, 17.7.1).

\numpara{8.4}
In the two statements which conclude this expos\'e, we return to the case of a field $k$ of arbitrary characteristic.

When $k$ is of characteristic $0$ (resp.\ $p>0$), let $n$ be an integer $\geq1$ (resp.\ an integer $\geq1$ prime to $p$); in both cases, we say simply that $n$ is prime to the characteristic of $k$. Moreover, if $G$ is a group scheme over $k$, we denote by $n_G:G\to G$ the morphism of $k$-schemes which sends an element $x$ of $G(T)$ to $x^n\in G(T)$, as $T$ ranges over $k$-schemes.

\begin{enonce}{Proposition}{}\label{VIIA.8.4}
Let $G$ be an algebraic group over a field $k$ and $n$ an integer prime to the characteristic of $k$. Then $n_G:G\to G$ is an \textup{\'etale} morphism.
\end{enonce}
% typo?: source does not assume G commutative in this global etaleness assertion; retained.
\nde{In the statement, ``\'etale at the origin'' has been changed to ``\'etale'', and the following sentence added.}
According to VI B 1.3, it suffices to show that $n_G$ is \'etale at the origin. Let $A$ be the local ring of $G$ at the origin and $I$ the maximal ideal of $A$. According to II 3.9.4, the map $\Lie(n_G):\Lie(G)\to\Lie(G)$ induced by $n_G$ is the homothety of ratio $n$. It is therefore an isomorphism, as is the endomorphism induced by $n_G$ on $I/I^2=\Lie(G)^*$. If $k$ is of characteristic $0$, $G$ is smooth over $k$ (VI B 1.6.1, see also VII B 3.3.1), so the canonical morphism $S(I/I^2)\to\gr_I(A)$ is an isomorphism, where $\gr_I(A)$ denotes the graded object associated to the $I$-adic filtration. It follows that $n_G$ induces an automorphism of $\gr_I(A)$, hence also of the completion $\widehat A$ of $A$, and therefore $n_G$ is \'etale at the origin (cf.\ EGA~IV$_4$, 17.6.3).

If the characteristic is $p>0$ and $G$ is of height $\leq1$, then $A$ is isomorphic to the quotient of the symmetric algebra of $\omega_{G/k}=I/I^2$ by the ideal generated by the $p$-th powers of elements of $\omega_{G/k}$ (cf.\ 7.4); one can then apply the ``same'' argument as in characteristic $0$, and one obtains that $n_G$ induces an automorphism of $A$.

If $G$ is of height $\leq r$ and we suppose our assertion proved for groups of height $\leq r-1$, denote by $B,A,A'$ the affine algebras of ${}_{\mathrm{Fr}}G$, $G$, and $G'={}_{\mathrm{Fr}}G\backslash G$, and by $n_B,n_A,n_{A'}$ the endomorphisms of $B,A,A'$ induced by $n_{{}_{\mathrm{Fr}}G}$, $n_G$, and $n_{G'}$.\nde{The original has been expanded in what follows.} Let $I'=I\cap A$ be the maximal ideal of $A'$; since there is a cartesian square
\[
\xymatrix@C=4em{{}_{\mathrm{Fr}}G\ar[r]\ar[d]&G\ar[d]\\e\ar[r]&G'}
\]
one has $B=A/I'A$. Observe that $n_{A'}$ (resp.\ $n_B$) is none other than the endomorphism induced by $n_A$ on $A'$ (resp.\ on $B$). According to VI A 3.2, $A$ is a faithfully flat $A'$-module, and since $A'$ is a local artinian ring ($G'$ being an algebraic $k$-group of height $\leq r-1$), it follows that $A$ is a free $A'$-module. Since the restriction of $n_A$ to $A'$ is $n_{A'}$, which is an isomorphism by the induction hypothesis, it follows from Nakayama's lemma that $n_A$ will be an automorphism if the endomorphism it induces on $A/I'A$ is one. Now this endomorphism is none other than $n_B$, which is an automorphism since $B$ is of height $\leq1$. Thus $n_A$ is an automorphism.
% typo?: source I'=I cap A (without prime on A), and B is said to have height; retained.

Finally, when $G$ is an arbitrary algebraic group over a field of characteristic $p>0$, what precedes shows that $n_G$ induces automorphisms of the $k$-schemes ${}_{\mathrm{Fr}^r}G$; these schemes are affine over $k$ and have as algebras the quotients of the local algebra $A$ by the ideal $I^{\{p^r\}}$ generated by the $p^r$-th powers of elements of $I$. Since $n_G$ defines automorphisms of the algebras $A/I^{\{p^r\}}$, one sees, by passage to the projective limit, that $n_G$ induces an automorphism of $\widehat A$, hence $n_G$ is \'etale at the origin (EGA~IV$_4$, 17.6.3).

\begin{proposition}{8.5}\label{VIIA.8.5}
Let $G$ be a finite algebraic group of rank $n$ over the field $k$. Then $n_G:G\to G$ is the zero morphism of $G$.
\end{proposition}
Let us point out at once the following corollary, obtained by combining 8.4 and 8.5:\nde{This corollary, indicated implicitly in the original by ``(confer 8.4)'', has been added. For another proof of the corollary, not using 8.5, see for example [TO70], Lemma 5.}
\begin{corollary}{8.5.1}\label{VIIA.8.5.1}
Let $G$ be a finite algebraic group of rank $n$ over the field $k$. If $n$ is prime to the characteristic of $k$, then $G$ is \'etale over $k$.
\end{corollary}

Let us now prove 8.5. Let $H$ be a normal subgroup of $G$ of rank $m$ over $k$. Denote by $\lambda:H\times G\to G$ the morphism induced by the multiplication of $G$. Then, with the notation of VI A 3.2, there is a cartesian square:
\[
\xymatrix@C=4em{H\times G\ar[r]^\lambda\ar[d]_{\mathrm{pr}_2}&G\ar[d]^\pi\\G\ar[r]_\pi&H\backslash G}
\]
Since $\pi:G\to H\backslash G$ is faithfully flat and quasi-compact (VI A 3.2), and $\mathrm{pr}_2$ is locally free of rank $m$, it follows from EGA~IV$_2$, 2.5.2 that $G\to H\backslash G$ is locally free of rank $m$. Denoting $\mathrm{rg}_k(H\backslash G)$ by $r$, one therefore has $n=\mathrm{rg}_k(G)=rm$.

On the other hand, one has an exact sequence of ``abstract'' groups
\[
1\longrightarrow H(T)\longrightarrow G(T)\longrightarrow(H\backslash G)(T)
\]
for every $k$-scheme $T$; it is therefore clear that $n_G$ is zero if $m_H$ and $r_{H\backslash G}$ are. Taking for $H$ the neutral component $G^0$ of $G$, $G^0\backslash G$ is \'etale (cf.\ VI A 5.5.1), so that one may suppose $G$ \'etale over $k$ or infinitesimal (cf.\ 7.0).

If $G$ is \'etale, one reduces, by extension of the base field, to the case where $k$ is algebraically closed. In this case, $G$ is a constant group (cf.\ I 4.1), and the statement is classical.

If $G$ is infinitesimal and nonzero, $k$ is necessarily of characteristic $p>0$ (cf.\ VI B 1.6.1 or VII B 3.3.1); the subgroups ${}_{\mathrm{Fr}^n}G$ then form a composition series of $G$ whose quotients are of height $\leq1$.

This reduces us to the case where $G$ is of height $\leq1$. Let $A$ (resp.\ $L$) be the affine algebra (resp.\ Lie algebra) of $G$, and $U=U_p(L)$. According to 7.4, one has $G=G_p(L)$, whence $A=U^*$; thus, if $\dim_kL=r$, the rank of $G$ over $k$ is $p^r$ (cf.\ 5.3.3). We shall therefore study the morphism $p_G:G\to G$ defined by raising to the power $p$; it induces an endomorphism $p_A$ of $A$ and, by duality, an endomorphism $p_U$ of $U$.

Let $I$ be the augmentation ideal of $A$; we shall show that $p_A(I)\subset I^p$. Assuming this established, one will therefore have $p_A^r(I)\subset I^{p^r}$. On the other hand, one knows that $I^{r(p-1)+1}=0$ (since $I$ is generated by $r$ elements with zero $p$-th power). Since $p^r>r(p-1)$, it follows that $p_A^r(I)=0$, hence $p_G^r$ is the zero morphism. It therefore remains to show the assertion:
\begin{equation}\tag{$*$}p_A(I)\subset I^p.\end{equation}
For every integer $s\geq1$, we shall denote by $m_A^{s-1}:A^{\otimes s}\to A$ (resp.\ $\Delta_U^s:U\to U^{\otimes s}$) the map induced by the multiplication $m_A$ of $A$ (resp.\ the comultiplication $\Delta_U$ of $U$). Then $p_A$ is equal to the following composite:
\[
A\xrightarrow{\Delta_A^{p-1}}A^{\otimes p}\xrightarrow{m_A^{p-1}}A,
\]
and since the transpose of $m_A$ (resp.\ $\Delta_A$) is $\Delta_U$ (resp.\ $m_U$), the endomorphism $p_U={}^tp_A$ of $U$ is the composite below:
\[
U\xrightarrow{\Delta_U^{p-1}}U^{\otimes p}\xrightarrow{m_U^{p-1}}U.
\]
% typo?: source initially writes Delta_U^s but later Delta_U^{s-1}; retained.
\nde{The original has been expanded in the following paragraph.}
Let $J$ be the augmentation ideal of $U$; one has $U=k1_U\oplus J$, and we shall denote by $\pi$ the projection $U\to J$ with kernel $k1_U$. For every integer $s\geq1$, denote by $(I^s)^\perp$ the orthogonal of $I^s$ in $A^*=U$, \ie $(I^s)^\perp$ is the set of $u\in U$ such that the composite below is zero:
\[
I^{\otimes s}\xrightarrow{m_A^{s-1}}I\xrightarrow{u}k.
\]
Since the transpose of $m_A$ is $\Delta_U$, one sees that $(I^s)^\perp$ is the vector subspace $P_{s-1}$ consisting of the $u\in U$ such that $\Delta_U^{s-1}(u)$ vanishes on $I^{\otimes s}$, \ie denoting by $\overline{\Delta_U^{s-1}}$ the composite of $\Delta_U^{s-1}$ and the projection $\pi^{\otimes s}:U^{\otimes s}\to J^{\otimes s}$, one obtains that
\[
(I^s)^\perp=P_{s-1}=\Ker\overline{\Delta_U^{s-1}}
\]
(see also VII B, 1.3.6). Thus, to show assertion $(*)$, one must show that the transpose map $p_U={}^tp_A$ maps $P_{p-1}$ into $I^\perp=k1_U$. Since $p_U(1_U)=1_U$, it suffices to show the assertion below, where $P_{p-1}^+$ denotes $J\cap P_{p-1}$:
\begin{equation}\tag{$**$}p_U(P_{p-1}^+)=0.\end{equation}
On the other hand, one shows easily, by induction on $s$, that $P_{s-1}^+$ is the vector subspace of $U$ generated by the products $x_1\cdots x_t$, with $1\leq t\leq s-1$ and $x_i\in L$ (cf.\ VII B 4.3). Now, if $x_1,x_2,\ldots,x_t$ are elements of $L$, one has:
\[
p_U(x_1x_2\cdots x_t)=m_U^{p-1}\left(\prod_{j=1}^t\sum_{i=1}^p1\otimes\cdots\otimes\overset{\substack{i\\\downarrow}}{x_j}\otimes\cdots\otimes1\right).
\]
It is clear that the expression $\prod_j\sum_i1\otimes\cdots\otimes x_j\otimes\cdots\otimes1$ is a sum of $p^t$ terms $x_h$ indexed by the maps $h$ from $\{1,\ldots,t\}$ to $\{1,\ldots,p\}$.\nde{The original has been expanded in what follows, replacing the notion of preorder by the equivalent notion of ordered partition.} Such a map $h$ defines an \emph{ordered partition} $\mathfrak p_h$ of $\{1,\ldots,t\}$ into at most $p$ parts. Indeed, denote by $i_1<\cdots<i_r$ the elements of the image of $h$, and, for $s=1,\ldots,r$, put $I_s=h^{-1}(i_s)$ and $x_{I_s}=\prod_{j\in I_s}x_j$, the product being taken in increasing order. Then $h$ corresponds to the $p$-tensor
\[
1\otimes\cdots\otimes x_{I_1}\otimes\cdots\otimes x_{I_r}\otimes\cdots\otimes1
\]
(where each $x_{I_s}$ is in position $i_s$), and its image under $m_U^p$ is the product:
\[
x_{I_1}\otimes\cdots\otimes x_{I_r}
\]
which depends only on the ordered partition $\mathfrak p=(I_1,\ldots,I_r)$, and which will be denoted by $x_{\mathfrak p}$. For fixed $\mathfrak p$, $x_{\mathfrak p}$ is obtained for all choices of $i_1<\cdots<i_r$ in $\{1,\ldots,p\}$, and one therefore obtains the equality
\[
p_U(x_1x_2\cdots x_t)=\sum_{\mathfrak p}\binom p{n(\mathfrak p)}x_{\mathfrak p},
\]
where $\mathfrak p$ ranges over the set of ordered partitions of $\{1,\ldots,t\}$ into at most $p$ parts, and $n(\mathfrak p)$ denotes the number of parts of $\mathfrak p$. (One has $1\leq n(\mathfrak p)\leq\min(t,p)$.)
% typo?: source m_U^p and tensor signs in the displayed "product" retained.
When $t<p$, all the terms $\binom p{n(\mathfrak p)}$ are therefore zero, so that $p_U(x_1\cdots x_t)=0$. Thus $p_U$ vanishes on $P_{p-1}^+$, which proves assertion $(**)$, hence $(*)$, and completes the proof of 8.5.

\begin{corollary}{8.5.2}\label{VIIA.8.5.2}
\nde{This corollary, mentioned in Exp.\ VIII, Remark 7.3.1, has been added.}
Let $S$ be a reduced scheme and $G$ a finite locally free $S$-group of rank $n$. Then $n_G:G\to G$ is the zero morphism of $G$.
\end{corollary}
Indeed, let $S'$ be the sum of the $\Spec\OO_{S,\eta}$, as $\eta$ ranges over the maximal points of $S$. Since $S$ is reduced, the morphism $S'\to S$ is schematically dominant, and the same is true of the morphism $f:G_{S'}\to G$, since $G$ is finite locally free over $S$ (cf.\ EGA~IV$_3$, 11.10.5). Since $G\to S$ is affine, hence separated, the coincidence locus of $n_G$ and the zero morphism is a closed subscheme of $G$; now it contains $f$ according to 8.5, hence equals $G$, \ie $n_G$ is the zero morphism.
% typo: source "Comme G -> S affine" omits est; supplied in English.

\begin{remark}{8.5.3}\label{VIIA.8.5.3}
Let us also point out that, according to a theorem of P. Deligne (see [TO70], p.\ 4), if $G$ is a finite locally free \emph{commutative} $S$-group of rank $n$ over an arbitrary base $S$, then $n_G=0$.
\end{remark}

\begin{thebibliography}{DG70}
\bibitem[BAlg]{VIIA.BAlg} N. Bourbaki, \textit{Alg\`ebre}, Chap.\ I--III, Hermann, 1974, Chap.\ X, Masson, 1980.
\bibitem[BAC]{VIIA.BAC} N. Bourbaki, \textit{Alg\`ebre commutative}, Chap.\ I--IV, Masson, 1985.
\bibitem[BLie]{VIIA.BLie} N. Bourbaki, \textit{Groupes et alg\`ebres de Lie}, Chap.\ I, Hermann, 1971.
\bibitem[DG70]{VIIA.DG70} M. Demazure, P. Gabriel, \textit{Groupes alg\'ebriques}, Masson \& North-Holland, 1970.
\bibitem[Ja03]{VIIA.Ja03} J. C. Jantzen, \textit{Representations of algebraic groups}, Academic Press, 1987; 2nd edition, Amer.\ Math.\ Soc., 2003.
\bibitem[TO70]{VIIA.TO70} J. Tate, F. Oort, \textit{Group schemes of prime order}, \textit{Ann. scient. \'Ec. Norm. Sup.} \textbf{3} (1970), 1--21.
% typo: source "Groups schemes of prime order" -> "Group schemes of prime order".
\end{thebibliography}
